\chapter{环与算术}\label{ch:b3:rings}

整数可唯一地分解为质数之积；多项式
在域上亦然。这两个事实是同一个定理吗？本章对此
作答：是的，并找出使交换环中的“算术”成为可能的
确切条件，即如下的链条：
\[
\text{欧几里得整环} \;\Longrightarrow\; \text{主理想整环}
\;\Longrightarrow\; \text{唯一分解整环 (UFD)},
\]
所有蕴含关系均已证明，所有逆命题均已被否定。该理论
随后在其大显身手之处接受检验：高斯整数
$\Z[\iu]$（它将在
周末问题中破解费马平方和定理），多元多项式环（高斯
引理、艾森斯坦判别法），以及诺特环，
并以希尔伯特基定理作结。全文中，\emph{环}指交换环，
且带单位元 $1 \neq 0$；本科二年级分册中关于 $\Z$ 与 $K[X]$ 的理想，
是我们的两个指导性例子。

\section{理想、商与同构定理}

\begin{definition}\label{def:b3:rings:ideal}
$I$ 为环 $A$ 的\emph{理想}\index{理想}，是满足
$AI \subseteq I$ 的加法子群。\emph{商环}\index{商环} $A/I$ 是商群
$(A, +)/I$，其乘法为
$(a + I)(b + I) = ab + I$，此乘法良定义，
因为把 $a$ 换为 $a + x$、$b$ 换为 $b + y$（$x, y \in I$）时，
$ab$ 的改变量为 $ay + xb + xy \in I$。投影 $\pi \colon A
\to A/I$ 是满的环同态，核为 $I$，且环同态的核
恰为全部理想。
\end{definition}

\begin{theorem}[第一同构定理]\label{thm:b3:rings:firstiso}
若 $f \colon A \to B$ 为环同态，则 $\bar f\colon
A/\ker f \to \operatorname{im} f$，$a + \ker f \mapsto f(a)$，是一个
环同构。更一般地，$f$ 经 $A/I$ 分解，对任一
理想 $I \subseteq \ker f$。$A/I$ 的理想恰为
$J/I$，其中 $J \supseteq I$ 为 $A$ 的理想（对应定理）。
\end{theorem}

\begin{proof}
与群的情形相同（见\cref{thm:b3:groups:firstiso,%
thm:b3:groups:correspondence}），注意所涉映射亦保持乘法：
$\bar f$ 良定义，为到像上的双射，且保持乘法；
对应 $J \mapsto J/I$，
$\bar J \mapsto \pi^{-1}(\bar J)$ 在两个方向上均保持理想，
因为 $\pi$ 为满的环同态。
\end{proof}

\begin{definition}\label{def:b3:rings:primemaximal}
设 $I \subsetneq A$ 为真理想。$I$ 称为
\emph{素理想}\index{素理想}，若 $ab \in I \Rightarrow a \in I$
或 $b \in I$；$I$ 称为\emph{极大理想}\index{极大理想}，若
不存在严格位于 $I$ 与 $A$ 之间的理想。
\end{definition}

\begin{proposition}\label{prop:b3:rings:primemaximal}
$I$ 为素理想 $\iff$ $A/I$ 为整环；$I$ 为极大理想
$\iff$ $A/I$ 为域。特别地，极大理想都是素理想。
\end{proposition}

\begin{proof}
记 $\bar a$ 为 $A/I$ 中的剩余类。“$I$ 为素理想”逐字译为
“$\bar a\bar b = 0 \Rightarrow \bar a = 0$ 或 $\bar b
= 0$”，而 $A/I \neq 0$ 译为 $I \neq A$，这正是
整环的定义。对于极大性，利用对应定理：不存在
严格位于 $I$ 与 $A$ 之间的理想 $\iff$ $A/I$ 除
$0$ 与自身外不再有其他理想 $\iff$ $A/I$ 为域 --- 最后一步的理由如下：在
域中，唯一的理想为 $0$ 与全环（若理想包含
某个 $x \ne 0$，则必含 $x^{-1}x = 1$）；反之，若每个非零 $x$
生成单位理想，则 $xy = 1$ 对某个 $y$ 成立。域都是
整环，故极大理想为素理想。
\end{proof}

\begin{example}\label{ex:b3:rings:primemaximal}
在 $\Z$ 中：素理想为 $(0)$ 与 $(p)$，$p$ 为质数；
极大理想为 $(p)$（$\Z/p\Z = \mathbb F_p$ 为域，
$\Z/(0) = \Z$ 则不是）。在 $K[X, Y]$ 中：$(X) \subsetneq (X, Y)$ 均
为素理想（$K[X,Y]/(X) \cong K[Y]$，为整环；
$K[X,Y]/(X,Y) \cong K$，为域），故 $(X)$ 是素理想而
非极大理想。
\end{example}

为保证极大理想在一般情形下\emph{存在}，
我们需要一条集合论原理。称偏序集
\emph{归纳的}，若其每个全序子集（\emph{链}）
均有上界。

\begin{theorem}[佐恩引理]\label{thm:b3:rings:zorn}
每个非空的归纳偏序集都有一个
极大元。\index{佐恩引理}
\admitted
\end{theorem}

\begin{remark}\label{rem:b3:rings:zorn}
这并非普通数学中的定理，而是一条\emph{公理}：
在集合论的策梅洛–弗兰克尔基本公理之上，它与
选择公理（“每族非空集合之积非空”）等价，
本书全文承认之。我们对每处用到它的地方都加以标注。
分析学中还将再次用到它（哈恩–巴拿赫定理，
见\cref{ch:b3:banach}）。
\end{remark}

\begin{theorem}[克鲁尔定理]\label{thm:b3:rings:krull}
每个真理想 $I \subsetneq A$ 都包含于某个
极大理想。\index{克鲁尔定理}
\end{theorem}

\begin{proof}
按包含关系排序集合 $\mathcal E$，它是所有
包含 $I$ 的真理想所成之集，它非空（$I \in \mathcal E$）。一条链
$(J_\lambda)$ 在 $\mathcal E$ 中有上界 $J = \bigcup
J_\lambda$：该并为理想（任意 $a, b \in J$ 同属某个 $J_\lambda$，
由全序性），为真理想（$1 \notin J_\lambda$ 对一切 $\lambda$ 成立），
且包含 $I$。由佐恩引理，$\mathcal
E$ 有极大元，它即为包含 $I$ 的极大理想（严格包含它的
真理想必属于 $\mathcal E$）。
\end{proof}

\begin{theorem}[中国剩余定理]\label{thm:b3:rings:crt}
设 $I_1, \dots, I_n$ 为 $A$ 中两两\emph{互素理想}\index{互素理想
}（$I_k + I_l = A$，$k \neq l$）。则
\[
A\Big/\bigcap_{k=1}^n I_k \;\xrightarrow{\;\sim\;}\;
\prod_{k=1}^n A/I_k,
\qquad
a \longmapsto (a + I_1, \dots, a + I_n),
\]
此外 $\bigcap_k I_k = I_1 I_2 \cdots I_n$（由乘积
生成的理想）。\index{中国剩余定理}
\end{theorem}

\begin{proof}
映射 $f(a) = (a + I_k)_k$ 是环同态，其核为
$\bigcap I_k$；由\cref{thm:b3:rings:firstiso}，只需
证明满射。固定 $k$；对每个 $l \neq k$，写 $1 = u_l +
v_l$，其中 $u_l \in I_k$，$v_l \in I_l$（由互素性）。于是
\[
e_k = \prod_{l \neq k} v_l = \prod_{l\neq k}(1 - u_l)
\equiv 1 \pmod{I_k},
\qquad e_k \in I_l \ (l \neq k),
\]
故 $f(e_k) = (0, \dots, 1, \dots, 0)$；对给定的目标 $(a_k +
I_k)_k$，元素 $\sum_k a_k e_k$ 的像恰为它。

乘积与交集的关系：$I_1\cdots I_n \subseteq \bigcap I_k$
恒成立。反之，由归纳法，只需处理 $n = 2$
（验证 $I_1$ 与 $I_2\cdots I_n$ 互素：将
各式 $1 = u_l + v_l$（$l \geq 2$）相乘，得 $1 \in I_1 + I_2\cdots
I_n$）。当 $n = 2$ 时，写 $1 = u + v$，$u \in I_1$，$v \in I_2$；
对 $x \in I_1 \cap I_2$，有 $x = xu + xv \in I_2I_1 + I_1I_2 =
I_1I_2$。
\end{proof}

\begin{example}\label{ex:b3:rings:crtz}
在 $\Z$ 中，设 $I_k = (m_k)$，$m_k$ 两两互素，则 $\Z/(m_1\cdots
m_n)\Z \cong \prod \Z/m_k\Z$ --- 即本科二年级分册中的中国剩余
定理。限制到单位群：$(\Z/mn\Z)^\times \cong
(\Z/m\Z)^\times \times (\Z/n\Z)^\times$，其中 $\gcd(m,n)=1$，由此
可得欧拉 $\varphi$ 函数的积性
（见\cref{exo:b3:rings:8}）。
\end{example}

\section{整除性：欧几里得、主理想、唯一分解}

\begin{definition}\label{def:b3:rings:divisibility}
设 $A$ 为整环，$a, b \in A$。称 $a$
\emph{整除} $b$（$a \mid b$），若 $b \in (a) = aA$。元素 $a,
b$ 称为\emph{相伴}\index{相伴}，若 $a = ub$，
其中 $u \in A^\times$（等价地 $(a) = (b)$）。非零
非单位元 $p$ 称为：
\begin{itemize}
  \item \emph{不可约元}\index{不可约元}，若 $p = ab$
        必有 $a \in A^\times$ 或 $b \in A^\times$；
  \item \emph{素元}\index{素元}，若 $p \mid ab$，必有 $p
        \mid a$ 或 $p \mid b$（即理想 $(p)$ 是素理想）。
\end{itemize}
\end{definition}

\begin{proposition}\label{prop:b3:rings:primeirred}
在任意整环中，素元 $\Rightarrow$ 不可约元。逆命题一般
不成立：在 $\Z[\iu\sqrt 5] = \{a + \iu b\sqrt5 : a, b
\in \Z\}$ 中，元素 $2$ 为不可约元，但不是素元。
\end{proposition}

\begin{proof}
设 $p$ 为素元，且 $p = ab$。则 $p \mid ab$，不妨设 $p \mid a$：
$a = pc$，于是 $p = pcb$，消去 $p$（整环中可消去）得 $cb =
1$，故 $b \in A^\times$。

在 $\Z[\iu\sqrt5]$ 中，取范数 $N(x + \iu y\sqrt 5) = x^2 +
5y^2$，它具有积性（即 $\abs z^2$）。若 $2 = ab$
且 $a, b$ 均为非单位元，则 $4 = N(a)N(b)$，其中 $N(a), N(b) \neq
1$（范数为 $1$ 的元素恰为 $\pm1$，即单位），故 $N(a) = 2$：
不可能，因为 $x^2 + 5y^2 = 2$ 无整数解。故 $2$ 为
不可约元。但 $2 \mid 6 = (1 + \iu\sqrt5)(1 - \iu\sqrt5)$
而 $2$ 不整除任一因子（$\frac12 \pm \frac{\iu\sqrt5}2
\notin \Z[\iu\sqrt5]$），故不是素元。
\end{proof}

\begin{definition}\label{def:b3:rings:pidufd}
称整环 $A$ 为：
\begin{itemize}
  \item \emph{欧几里得整环}\index{欧几里得整环}，若存在映射
        $\nu \colon A \setminus \{0\} \to \N$（称为\emph{欧几里得函数
        }）使得对所有 $a, b$（$b \ne 0$），存在
        $q, r$ 使 $a = bq + r$，且（$r = 0$ 或 $\nu(r) <
        \nu(b)$）；
  \item \emph{主理想}（\emph{PID}\index{主理想整环
        }）若每个理想都形如 $(a)$；
  \item \emph{唯一分解}（\emph{UFD}\index{唯一分解整环
        }），即每个非零非单位元都是
        不可约元之积，且在次序及相伴意义下唯一。
\end{itemize}
\end{definition}

\begin{theorem}\label{thm:b3:rings:euclideanpid}
欧几里得整环 $\Rightarrow$ 主理想整环。
\end{theorem}

\begin{proof}
设 $I \neq (0)$ 为理想，$b \in I \setminus\{0\}$ 使得
$\nu(b)$ 最小。对 $a \in I$ 做除法：$a = bq + r$；则 $r = a
- bq \in I$，若 $\nu(r) < \nu(b)$，则与极小性矛盾，故
$r = 0$，从而 $a \in (b)$，即 $I = (b)$。
\end{proof}

\begin{example}\label{ex:b3:rings:euclidean}
$\Z$（取 $\nu = \abs\cdot$）与 $K[X]$（取 $\nu = \deg$）都是
欧几里得整环 --- 本科二年级分册已证明过这两种除法。同样，
$\Z[\iu]$ 也是，取 $\nu = N$ 为平方范数
（见\cref{exo:b3:rings:4}）；证明的几何直观见
下图。存在不是欧几里得整环的主理想整环，但验证颇为精细
（标准例子是
$\Z\bigl[\frac{1+\iu\sqrt{19}}2\bigr]$）；不是主理想整环的唯一分解整环
很容易给出：$K[X, Y]$（见\cref{exo:b3:rings:6}），或 $\Z[X]$。
\end{example}

\begin{omfigure}
\begin{tikzpicture}[scale=1.05]
  \clip (-2.7,-1.7) rectangle (2.7,2.2);
  \foreach \x in {-2,...,2}
    \foreach \y in {-1,...,2}
      \fill[black!55] (\x,\y) circle (1.4pt);
  \draw[omDef, thick] (0.6,0.4) circle (1);
  \fill[omThm] (0.6,0.4) circle (1.8pt)
    node[below right, font=\small, omThm] {$a/b$};
  \fill[omDef] (1,0) circle (2pt)
    node[below, font=\small, omDef] {$q$};
  \draw[omDef, ->, shorten >=2.5pt, shorten <=2.5pt]
    (0.6,0.4) -- (1,0);
\end{tikzpicture}

{\small 高斯整数中的除法：精确商
$a/b \in \C$ 与某格点的距离 $\leq \frac{\sqrt2}{2} < 1$，
其中 $q \in \Z[\iu]$；于是 $r = a - bq$ 满足 $N(r) =
N(b)\,\abs{a/b - q}^2 < N(b)$。一次欧几里得除法，便足以
推出整套算术。}
\end{omfigure}

\begin{lemma}[主理想的升链]
\label{lem:b3:rings:acc}
在主理想整环中，每个递增的理想序列 $I_1 \subseteq I_2
\subseteq \cdots$ 必从某项起恒定。
\end{lemma}

\begin{proof}
$I = \bigcup_n I_n$ 是理想（并集递增），因此 $I =
(a)$；元素 $a$ 属于某个 $I_N$，于是 $I = (a)
\subseteq I_N \subseteq I_n \subseteq I$，对 $n \geq N$ 成立。
\end{proof}

\begin{lemma}[贝祖；欧几里得引理]\label{lem:b3:rings:bezout}
设 $A$ 为主理想整环，$a, b \in A$。则 $(a) + (b) = (d)$，
其中 $d$ 为一个\emph{最大公因数}\index{最大公因数}：
$d \mid a$，$d \mid b$，且 $a,
b$ 的任一公因数都整除 $d$；并有 $d = au + bv$，其中 $u, v$ 为某些元素，
（贝祖）。于是主理想整环中每个不可约元
都是素元。
\end{lemma}

\begin{proof}
$(a) + (b)$ 是理想，因而为 $(d)$；由 $a, b \in (d)$ 得 $d \mid
a, b$；并且 $d = au + bv \in (a) + (b)$。设 $c$ 为
$a, b$ 的公因数，则其整除 $au + bv = d$。

欧几里得：设 $p$ 为不可约元，$p \mid ab$，$p \nmid a$。
取 $d$ 为 $p$ 与 $a$ 的一个最大公因数，则它整除 $p$，于是 $d$ 为单位或
与 $p$ 相伴（由不可约性）；由 $p \nmid a$ 排除相伴。
故 $1 = pu + av$，从而 $b = pub + abv$，且 $p$ 整除两项：
$p \mid b$。
\end{proof}

\begin{theorem}\label{thm:b3:rings:pidufd}
主理想整环 $\Rightarrow$ 唯一分解整环。
\end{theorem}

\begin{proof}
\emph{存在性。}假设某个非零非单位元 $a$ 不能分解为
不可约元之积。则 $a$ 不是不可约元：$a =
a_1b_1$，两因子均非单位；其中至少有一个，不妨设为
$a_1$，仍没有分解（两个可分解元之积
仍可分解）。迭代可得 $a = a_0, a_1, a_2,
\dots$，每个都是后者的真因子且无分解，
故 $(a_0) \subsetneq (a_1) \subsetneq (a_2) \subsetneq \cdots$
——包含是严格的，因为在 $a_n = a_{n+1}c$ 中 $c$ 为
非单位；若 $(a_n) = (a_{n+1})$，则将迫使 $c \in A^\times$（整环中可消去）。
这与\cref{lem:b3:rings:acc}矛盾。

\emph{唯一性。}设 $p_1 \cdots p_r = q_1 \cdots q_s$，
其中因子均不可约，$r \leq s$，对 $r$ 归纳。
由\cref{lem:b3:rings:bezout}，素元 $p_1$ 整除右端，
于是它整除某个 $q_j$；不妨设 $j = 1$。由于 $q_1$ 不可约，
且 $p_1$ 非单位，故 $q_1 = u p_1$，其中 $u \in A^\times$：$p_1,
q_1$ 相伴。消去 $p_1$：$p_2 \cdots p_r = (u q_2)
q_3\cdots q_s$，再由归纳法即得（$r = 1$ 时须 $s = 1$：
单位乘以不可约元不可能等于 $1$）。
\end{proof}

\begin{remark}\label{rem:b3:rings:ufdgcd}
在唯一分解整环中，最大公因数存在（取分解式中指数的较小者），
且欧几里得引理成立——不可约元 $=$ 素元
（见\cref{exo:b3:rings:2}）——但贝祖恒等式可能不成立：在 $\Z[X]$ 中，
$\gcd(2, X) = 1$，但 $1 \neq 2U + XV$（代入 $X = 0$：
$1 = 2U(0)$，不可能）。贝祖恒等式是
主理想整环所特有的性质。
\end{remark}

\begin{example}[不具唯一分解的环]
\label{ex:b3:rings:nonufd}
这些蕴含都不是等价的：欧几里得整环 $\Rightarrow$ 主理想整环
$\Rightarrow$ 唯一分解整环；其中最后一个的失效
值得完整地看一次：
\[
A = \Z[\iu\sqrt5] = \{a + \iu b\sqrt5 : a, b \in \Z\},
\qquad N(a + \iu b\sqrt5) = a^2 + 5b^2,
\]
在此环中，范数是可乘的，且 $N(z) = 1$ 当且仅当 $z \in
A^\times = \{\pm1\}$。考虑
\[
6 = 2 \cdot 3 = (1 + \iu\sqrt5)(1 - \iu\sqrt5).
\]
四个因子都不可约：它们的范数分别为 $4, 9, 6,
6$，而真分解 $z = z_1z_2$ 将迫使
$N(z_1) \in \{2, 3\}$——但 $a^2 + 5b^2$ 永不等于 $2$
或 $3$（$b = 0$ 时只余非平方数 $2, 3$；$\abs b \geq
1$ 时给出 $\geq 5$）。然而 $2$ 与 $1 \pm
\iu\sqrt5$ 都不相伴（范数 $4 \neq 6$）：这是
$6$ 的两种本质不同的不可约分解。等价地，
不可约元 $\neq$ 素元：这里 $2$ 整除乘积 $(1 +
\iu\sqrt5)(1 - \iu\sqrt5) = 6$，但两个因子均不被它整除（同样可由范数看出）。
针对这一失效的理想论修补——
即对\emph{理想}而非元素作分解，这正是代数数论的诞生；
在本阶段，此例说明
本章的欧几里得整环 $\Z$、$K[X]$、
$\Z[\iu]$ 究竟有多特殊。
\end{example}

\begin{method}\label{met:b3:rings:quotient}
要识别商环 $A/I$，只需寻找满同态，
$f \colon A \to B$，其核为 $I$，并援引
\cref{thm:b3:rings:firstiso}；当 $A = C[X]$ 为多项式环时，
$f$ 通常取为求值映射。于是 $\Z[X]/(X^2+1) \cong
\Z[\iu]$（在 $\iu$ 处求值），$K[X,Y]/(Y - X^2) \cong K[X]$
（将 $Y$ 代入 $X^2$），$\R[X]/(X^2+1) \cong \C$。
要证明 $I$ 为素理想或极大理想，只需证明商环是整环或域
（\cref{prop:b3:rings:primemaximal}）。
\end{method}

\section{唯一分解整环上的多项式：高斯与艾森斯坦}

本节中 $A$ 为唯一分解整环，其分式域为 $K$
（以形式商 $a/b$（$b \neq 0$）构成的域，
与 $\Q$ 由 $\Z$ 得出的方式完全相同；本科二年级分册
已对 $\Q$ 作过这一构造，此构造可原样照搬）。本节的目标是
唯一分解性由 $A$ 传递到 $A[X]$，
而 $A$ 上的不可约性本质上就是更大的域 $K$ 上的不可约性。

\begin{definition}\label{def:b3:rings:content}
\emph{容度}\index{容度} $c(P)$ 是非零多项式 $P \in A[X]$ 的
系数的一个最大公因数（至多相差一个单位）；$P$ 称为\emph{本原}的，
若 $c(P) \in A^\times$。每个 $P
\in A[X]$ 都可写成 $P = c(P)\,P_1$，其中 $P_1$ 为本原的，
而每个 $P \in K[X]\setminus\{0\}$ 都可写成 $P = \lambda P_1$，其中 $\lambda
\in K^\times$ 且 $P_1 \in A[X]$ 为本原（先去分母，
再提取容度）。
\end{definition}

\begin{lemma}[高斯]\label{lem:b3:rings:gauss}
$A[X]$ 中两个本原多项式之积仍是本原的；
因此 $c(PQ) = c(P)c(Q)$（至多相差单位）。\index{高斯引理}
\end{lemma}

\begin{proof}
设 $P, Q$ 为本原多项式，并设某个不可约元（在唯一分解整环中即素元）
$p$ 整除 $PQ$ 的全部系数。模 $p$ 约化，
在 $(A/(p))[X]$ 中 $\bar P \bar Q = 0$。但 $A/(p)$ 是整环
（$(p)$ 为素理想），故 $(A/(p))[X]$ 是整环（首项系数相乘），
从而 $\bar P = 0$ 或 $\bar Q = 0$：$p$ 整除 $P$ 的全部系数或 $Q$ 的全部系数，
与本原性矛盾。至于后一结论，
设 $P = c(P)P_1$，$Q = c(Q)Q_1$，则 $PQ =
c(P)c(Q) P_1Q_1$，且 $P_1Q_1$ 为本原的。
\end{proof}

\begin{theorem}\label{thm:b3:rings:gaussufd}
设 $A$ 为唯一分解整环，其分式域为 $K$。
\begin{enumerate}
  \item 若 $P \in A[X]$ 为次数 $\geq 1$ 的本原多项式，
        则它在 $A[X]$ 中不可约，当且仅当它在 $K[X]$ 中不可约。
  \item $A[X]$ 为唯一分解整环；其不可约元即为 $A$ 的不可约元，
        以及在 $K$ 上不可约的本原多项式。特别地，
        $\Z[X]$，以及经归纳得到的 $K[X_1, \dots, X_n]$，
        和 $\Z[X_1, \dots, X_n]$ 均为唯一分解整环。
\end{enumerate}
\end{theorem}

\begin{proof}
(1)（$\Leftarrow$）若 $P = QR$ 为 $A[X]$ 中的分解，其中 $Q, R$ 均非单位，
则两个因子都不是常数（本原多项式的常数因子
必为单位），故该分解在 $K[X]$ 中仍是非平凡分解。
（$\Rightarrow$）设 $P = QR$，其中 $Q, R \in K[X]$，
次数 $\geq 1$。写 $Q = \lambda Q_1$，$R = \mu R_1$，其中 $Q_1, R_1
\in A[X]$ 为本原多项式，于是 $P = \lambda\mu\, Q_1R_1$，
而由高斯引理，$Q_1R_1$ 仍为本原的。取容度，得 $\lambda\mu \in A^\times$
（两边的容度都是单位；形式上，$\lambda\mu = c(P) \in
A^\times$ 至多相差单位，特别地 $\lambda \mu \in A$）：$P
= (\lambda\mu Q_1) R_1$ 是 $A[X]$ 中的非平凡分解。

(2) 存在性：设 $P \neq 0$ 且非单位，分解 $P = c(P)P_1$，
将 $c(P)$ 分解为 $A$ 中的不可约元之积，并
把 $P_1$ 在唯一分解整环 $K[X]$ 中写成 $\prod Q_i$，其中 $Q_i \in K[X]$ 不可约；
记 $Q_i = \lambda_i R_i$，其中 $R_i \in A[X]$ 为本原多项式（从而
在 $K$ 上不可约，由 (1) 知亦在 $A[X]$ 中不可约），
则乘积 $\prod \lambda_i$ 与前面一样是 $A$ 的单位，且 $P_1 = u\prod
R_i$。唯一性：比较一个分解的常数部分与
多项式部分；常数之积为 $c(P)$（由高斯引理），唯一性
由 $A$ 的唯一分解性保证；多项式部分
给出 $K[X]$ 中同一多项式的两个分解，因此二者
在相差 $K^\times$ 中常数的意义下对应（由 $K[X]$ 的唯一分解性，
见\cref{thm:b3:rings:pidufd}），而对应的本原多项式
在 $K[X]$ 中相伴，则在 $A[X]$ 中亦相伴：
若 $R = \lambda R'$，其中 $R, R'$ 为本原的，$\lambda \in K^\times$，则取容度
可得 $\lambda \in A^\times$。
\end{proof}

\begin{theorem}[不可约性判别法]\label{thm:b3:rings:criteria}
设 $A$ 为唯一分解整环，$K$ 为其分式域，且 $P = a_nX^n + \dots
+ a_0 \in A[X]$ 为次数 $n \geq 1$ 的本原多项式。
\begin{enumerate}
  \item （\emph{约化}）若 $p \in A$ 为素元，$p \nmid a_n$，
        且约化 $\bar P$ 在 $(A/(p))[X]$ 中
        不可约，则 $P$ 在 $K[X]$ 中不可约（从而
        在 $A[X]$ 中亦不可约）。
  \item （\emph{艾森斯坦判别法}\index{艾森斯坦判别法}）若存在
        某素元 $p$ 满足 $p \nmid a_n$，且 $p \mid a_i$（$0
        \leq i < n$），
        $p^2 \nmid a_0$，则 $P$ 在 $K[X]$ 中不可约（从而在 $A[X]$ 中亦然）。
\end{enumerate}
\end{theorem}

\begin{proof}
由\cref{thm:b3:rings:gaussufd}(1)，$K$ 上的非平凡分解
给出 $P = QR$，其中 $Q, R \in A[X]$，$\deg Q, \deg R \geq 1$
（常数被排除：它们要么是单位，要么破坏
本原性）。

(1) 模 $p$ 约化：$\bar P = \bar Q\bar R$ 在 $(A/(p))[X]$ 中成立。
由于 $p \nmid a_n$，且约化只会使 $\deg$ 降低，$\deg \bar
Q = \deg Q \geq 1$，$\deg\bar R = \deg R \geq 1$（它们的首项系数之积为 $\bar a_n \neq 0$，
因此两者次数均不降低）：
$\bar P$ 有非平凡分解 --- 矛盾。

(2) 模 $p$ 约化：$\bar Q \bar R = \bar P = \bar a_n X^n$（低次
系数全部消失）。在整环 $(A/(p))[X]$ 中，
$cX^n$（$c \ne 0$）的因式分解必为常数与纯
幂 $c'X^k$ 之积：事实上，若 $\bar Q\bar R = \bar a_nX^n$，且设
$\bar Q$ 在次数 $< \deg\bar Q$ 处有非零系数，取
最低次非零项：$\operatorname{val}(\bar Q\bar R) =
\operatorname{val}\bar Q + \operatorname{val}\bar R$（整环中），
它必等于 $n = \deg\bar Q + \deg\bar R$，于是
对二者均有 $\operatorname{val} = \deg$，即二者皆为单项式。如
上所述，约化不降次，故 $Q$ 与 $R$ 的常数项
$Q(0), R(0)$ 均被 $p$ 整除——两者皆然，
因两个约化式皆为次数 $\geq 1$ 的单项式。于是 $p^2 \mid
Q(0)R(0) = a_0$，矛盾。
\end{proof}

\begin{example}\label{ex:b3:rings:eisenstein}
$X^n - p$ 在 $\Q$ 上对每个质数 $p$ 及每个 $n
\geq 1$ 不可约（$p$ 处的艾森斯坦判别法）：$\Q$ 上
存在各次数的不可约多项式 --- 这与 $\C$（次数
$1$，达朗贝尔—高斯定理，在\cref{ch:b3:holomorphic}中证明）及
$\R$（次数 $1, 2$）形成鲜明对比。\emph{平移}的技巧
扩大了艾森斯坦判别法的适用范围：第 $p$ 个\emph{分圆
多项式}\index{分圆多项式} $\Phi_p = X^{p-1} +
\dots + X + 1 = \frac{X^p - 1}{X - 1}$，并有
\[
\Phi_p(X + 1) = \frac{(X+1)^p - 1}{X}
= X^{p-1} + \binom{p}{1}X^{p-2} + \dots + \binom{p}{p-1},
\]
由 $p$ 处的艾森斯坦判别法（$p \mid \binom pk$ 对 $0 < k < p$ 成立，且
$\binom{p}{p-1} = p \not\equiv 0 \bmod p^2$），$\Phi_p(X+1)$ 不可约，
从而 $\Phi_p$ 在 $\Q$ 上不可约。这正是
在\cref{ch:b3:galois}中讲述的正 $17$ 边形故事的代数核心。
\end{example}

\begin{method}\label{met:b3:rings:irreducibility}
要证明 $P \in \Z[X]$ 在 $\Q$ 上不可约，可依次：(i) 使 $P$
为本原多项式；(ii) 试用艾森斯坦判别法于 $P(X)$ 及其平移 $P(X \pm
1)$；(iii) 试对不整除首项系数的小素数
取模约化：只需模 \emph{某一个} $p$ 不可约即可，
而在 $\mathbb F_p$ 上不可约性可由有限次检验判定
（无根即排除次数为 $1$ 的因式；再检验有限多个
各次数 $\leq \deg P/2$ 的因式）；(iv) 若其余方法皆失效，
则用待定系数法。需注意：对每个 $p$ 模约化皆可约，
\emph{并不}蕴含 $\Q$ 上可约
（见\cref{exo:b3:rings:11}）。
\end{method}

\section{诺特环}

\begin{definition}\label{def:b3:rings:noetherian}
称环 $A$ 为\emph{诺特环}\index{诺特环}，若
$A$ 的每个理想都是有限生成的。
\end{definition}

\begin{proposition}\label{prop:b3:rings:noethacc}
$A$ 为诺特环当且仅当每个理想的递增序列
最终恒定（\emph{升链条件}），当且仅当每个
非空理想族（关于包含关系）皆有极大元。
\end{proposition}

\begin{proof}
\emph{（有限生成 $\Rightarrow$ 升链条件）}：设链 $I_1 \subseteq I_2
\subseteq \cdots$，其并 $I$ 为理想，由 $x_1,
\dots, x_r$ 生成；所有 $x_i$ 都属于某个 $I_N$，故 $I = I_N = I_n$，
对 $n \geq N$ 成立。\emph{（升链条件 $\Rightarrow$ 极大元）}：
若非空族 $\mathcal F$ 无极大元，取 $I_1 \in
\mathcal F$，再归纳地取 $I_{n+1} \supsetneq I_n$ 于
$\mathcal F$ 中（因 $I_n$ 非极大，故可行），得无限的
严格递增链。（此处用到依赖选择公理，这是一种较弱的选择公理，
此处不予深究。）
\emph{（极大元 $\Rightarrow$ 有限生成）}：给定理想 $I$，
由包含于 $I$ 的有限生成理想组成的族
非空（$(0)$ 属于其中）；其极大元 $J = (x_1, \dots, x_r)$ 必
等于 $I$：否则把 $x \in I \setminus J$ 加入
生成元，即得到该族中一个严格更大的成员。
\end{proof}

\begin{theorem}[希尔伯特基定理]\label{thm:b3:rings:hilbert}
若 $A$ 为诺特环，则 $A[X]$ 也是诺特环。由此 $A[X_1, \dots,
X_n]$ 及其任一商环亦然。\index{希尔伯特基定理}
\end{theorem}

\begin{proof}
设 $I$ 为 $A[X]$ 的理想，并假设 $I$ 不是有限
生成的。构造序列：取 $f_1 \in I \setminus \{0\}$ 为次数最小者，
再归纳地取 $f_{k+1} \in I \setminus (f_1,
\dots, f_k)$ 中次数最小者（该集合非空，依
假设）。次数 $d_k = \deg f_k$ 不减（由
各步选取的极小性：$f_{k+1}$ 在第
$k+1$ 步时可供选取——精确地说，$f_{k+1} \notin (f_1,\dots,f_k) \supseteq
(f_1, \dots, f_{k-1})$，故 $f_{k+1}$ 在第 $k$ 步已参与竞争，
结果落选或并列，故 $d_{k+1} \geq d_k$。设 $a_k \in A$ 为
$f_k$ 的首项系数。理想链 $(a_1) \subseteq (a_1,
a_2) \subseteq \cdots$ 稳定，即 $a_{n+1} \in (a_1, \dots,
a_n)$ 对某个 $n$ 成立，记 $a_{n+1} = \sum_{k\leq n} u_k a_k$。
考虑
\[
g = f_{n+1} - \sum_{k=1}^{n} u_k X^{\,d_{n+1} - d_k} f_k .
\]
于是 $g \in I \setminus (f_1, \dots, f_n)$（和式属于该理想，
而 $f_{n+1}$ 不属于），但
次数 $d_{n+1}$ 的系数相消，$\deg g < d_{n+1}$，与
$d_{n+1} = \deg f_{n+1}$ 的极小性矛盾。

反复应用，$A[X_1, \dots, X_n] = (A[X_1, \dots, X_{n-1}])[X_n]$
为诺特环；商环 $A/I$ 也是诺特环，因为它的理想
$J/I$ 可提升为 $A$ 的理想（由对应定理），其中
有限个生成元投影后仍为生成元。
\end{proof}

\begin{remark}\label{rem:b3:rings:noethmeaning}
诺特性是代数几何的有限性公理：任意
一组含 $n$ 个变量的多项式方程组，无论多么庞大，
都等价于其中有限个方程——其解集可由
有限个多项式截出。主理想整环（PID）是诺特环
（这是显然的）；含无穷多个变量的 $\Z[X_1, X_2, \dots]$
不是诺特环（$(X_1) \subsetneq (X_1, X_2) \subsetneq \cdots$）。
非诺特环在分析中也自然出现：$\intcc01$ 上的
连续函数构成一个非诺特环（见\cref{exo:b3:rings:10}）。
\end{remark}

\section{练习}

\begin{exercise}[$\star$]\label{exo:b3:rings:1}
确定下列商环：(a) $\Z[X]/(X^2 + 1) \cong \Z[\iu]$；
(b) $\R[X]/(X^2+1) \cong \C$；(c) $\mathbb F_2[X]/(X^2 + X + 1)$
是含 $4$ 个元素的域——写出其乘法表。
\end{exercise}

\begin{exercise}[$\star$]\label{exo:b3:rings:2}
(a) 证明：在唯一分解整环中，每个不可约元都是素元。
(b) 证明：有限整环是域。
(c) 由此推出：有限环中每个素理想都是极大理想。
\end{exercise}

\begin{exercise}[$\star$]\label{exo:b3:rings:3}
在 $\Z[\iu\sqrt5]$ 中：验证 $3$、$1 + \iu\sqrt5$ 与 $1 -
\iu\sqrt5$ 均为不可约元，且 $9 = 3\cdot 3 = (2 +
\iu\sqrt5)(2 - \iu\sqrt5)$，并由
\cref{prop:b3:rings:primeirred}再次得出 $\Z[\iu\sqrt5]$ 不是
唯一分解整环。唯一性究竟在哪里失效？
\end{exercise}

\begin{exercise}[$\star\star$]\label{exo:b3:rings:4}
(a) 证明：$\Z[\iu]$ 关于范数 $N(x + \iu y) =
x^2 + y^2$ 是欧几里得整环：给定 $a, b \neq 0$，取 $q \in \Z[\iu]$ 使其最接近
$a/b \in \C$。
(b) 求 $\Z[\iu]^\times$。
(c) 对 $\Z[\iu\sqrt2]$ 及 $N(x + \iu y\sqrt2) =
x^2 + 2y^2$ 讨论同样的问题。为何同样的论证对
$\Z[\iu\sqrt5]$ 不成立？
\end{exercise}

\begin{exercise}[$\star\star$]\label{exo:b3:rings:5}
设 $A$ 为环。(a) 证明：若 $x$ 幂零（即 $x^n = 0$
对某个 $n$），则 $1 + x \in A^\times$。(b) 证明：若 $A$ 为
整环，则 $A[X]^\times = A^\times$；在
$\Z/4\Z$ 上给出反例。(c) 证明：整环中幂等元（$e^2 = e$）只有
$0, 1$，幂零元只有 $0$。
\end{exercise}

\begin{exercise}[$\star\star$]\label{exo:b3:rings:6}
在 $A = K[X, Y]$ 中：(a) 证明理想 $(X, Y)$ 是极大理想但
不是主理想——因此 $K[X,Y]$ 是唯一分解整环
（\cref{thm:b3:rings:gaussufd}），却不是主理想整环；(b) 确定
$K[X, Y]/(Y - X^2)$ 与 $K[X,Y]/(XY - 1)$ 作为有理函数的
子环；(c) $(Y - X^2)$ 是素理想吗？是极大理想吗？
\end{exercise}

\begin{exercise}[$\star\star$]\label{exo:b3:rings:7}
判断下列多项式在 $\Q$ 上是否可约：$X^5 - 12X^3 + 36X - 12$；\quad
$X^4 + X + 1$ \emph{（模 $2$ 约化）}；\quad $X^4 + 4$；\quad
$\Phi_8 = X^4 + 1$ \emph{（作平移 $1$）}；\quad $X^3 - X - 1$。
\end{exercise}

\begin{exercise}[$\star\star$]\label{exo:b3:rings:8}
(a) 由\cref{thm:b3:rings:crt}证明：欧拉函数在互素的自变量上可乘，
且 $\varphi(p^k) =
p^{k-1}(p-1)$；由此推出 $\varphi(n) = n\prod_{p \mid n}(1 -
\frac1p)$。
(b) 解同余方程组：$x \equiv 2 \pmod 7$，$x \equiv 5 \pmod{11}$，$x
\equiv 1 \pmod{13}$，并给出
\cref{thm:b3:rings:crt}证明中的幂等元 $e_k$。
\end{exercise}

\begin{exercise}[$\star\star\star$]\label{exo:b3:rings:9}
（幂零根基）设 $\operatorname{Nil}(A)$ 为
幂零元之集。(a) 证明 $\operatorname{Nil}(A)$ 是
包含于每个素理想中的理想。(b) 反之，设 $a$ 为
非幂零元；利用对避开 $S =
\{a^n : n \in \N\}$ 的理想使用佐恩引理，构造一个不含 $a$ 的素理想。
由此得出：
\[
\operatorname{Nil}(A) = \bigcap_{\mathfrak p \text{ 素}}
\mathfrak p .
\]
\end{exercise}

\begin{exercise}[$\star\star\star$]\label{exo:b3:rings:10}
(a) 设 $A$ 为诺特环，$f \colon A \to A$ 为满
的环同态。证明 $f$ 是单射。\emph{（考虑 $\ker
f \subseteq \ker f^2 \subseteq \cdots$.)}
(b) 证明连续函数环 $\mathcal C(\intcc01, \R)$
不是诺特环。\emph{（考虑 $I_n = \{f : f = 0
\text{ 在 } \intcc0{1/n}\}$。）}
(c) 证明：在诺特\emph{整环}中，每个非零
非单位元都是（有限个）不可约元之积——
因此 $\Z[\iu\sqrt 5]$ 的非唯一分解性只是
\emph{唯一性}的失效。
\end{exercise}

\begin{exercise}[$\star\star\star$]\label{exo:b3:rings:11}
设 $P = X^4 + 1$。
(a) 证明 $P$ 在 $\Q$ 上不可约
（见\cref{exo:b3:rings:7}）。
(b) 证明：$P$ 模\emph{每个}素数 $p$ 都可约：
先处理 $p = 2$；再对奇素数 $p$，证明 $8 \mid p^2 - 1$，并
暂且承认（将在\cref{ch:b3:galois}中证明）：含 $p^2$ 个元素的域的
乘法群是循环群，由此推出 $P$ 模 $p$ 分裂为两个二次因式；
当 $-1$、$2$、$-2$ 之一是模
$p$ 的平方时，将这两个因式显式写出，
并证明其中之一总是成立。
\end{exercise}

\begin{exercise}[$\star\star$]\label{exo:b3:rings:12}
（幂等元分解环）设 $e$ 为交换环 $A$ 中的元素，
若 $e^2 = e$，则称之为\emph{幂等元}。
(a) 证明：若 $e$ 为幂等元，则 $1 - e$ 亦然，并且
映射 $x \mapsto (ex, (1-e)x)$ 是环同构 $A
\cong Ae \times A(1-e)$，其中 $Ae$ 是以 $e$ 为单位元的环。
(b) 求整环中的全部幂等元，以及 $\Z/12\Z$ 中的全部幂等元；
给出同构 $\Z/12\Z \cong \Z/4\Z \times \Z/3\Z$，
通过指出其两个非平凡幂等元。
(c) 证明 $\Z/n\Z$ 的中国剩余定理分解
（见\cref{ex:b3:rings:crtz}）恰好对应于
幂等元 $e_i \equiv 1 \bmod p_i^{a_i}$，$e_i \equiv 0$
模其余素数幂：环沿其幂等元分解，
正如空间沿投影分解。
\end{exercise}

\section{问题：费马两平方和定理}

\begin{problem}[{周末问题 --- 两平方和，借助
$\Z[\iu]$}]\label{pb:b3:rings:1}
哪些整数是两个平方数之和？费马的答案（1640 年）是
算术中的瑰宝之一；高斯整数把它的证明
化为环论。全文中，$N(x + \iu y) = x^2 + y^2$ 表示
范数，$\Z[\iu]$ 是欧几里得整环（见\cref{exo:b3:rings:4}），从而
是主理想整环与唯一分解整环，并称\emph{高斯素数}为
$\Z[\iu]$ 中的素元（即不可约元）。

\medskip
\noindent\textbf{第一部分 --- 范数与高斯素数。}
\begin{enumerate}
  \item 验证 $N(zw) = N(z)N(w)$，并再次推出 $\Z[\iu]^\times =
        \{\pm1, \pm\iu\}$，并证明\emph{婆罗摩笈多
        恒等式}：两个两平方和之积
        仍是两平方和。
  \item 证明：若 $N(z)$ 是质数，则 $z$ 是一个
        高斯素数。
  \item 证明：每个高斯素数 $\pi$ 恰整除一个
        质数 $p$ \emph{（考虑 $N(\pi) = \pi\bar\pi$）}，
        并且此时 $N(\pi) \in \{p, p^2\}$。
  \item 推出二分性：对每个质数 $p$，或者 $p$
        在 $\Z[\iu]$ 中仍为素元（且不存在范数为
        $p$ 的高斯素数），或者 $p = \pi\bar\pi$，其中 $\pi$ 为
        范数为 $p$ 的高斯素数 --- 此时 $p = a^2 + b^2$。
\end{enumerate}

\medskip
\noindent\textbf{第二部分 --- 威尔逊定理与 $-1$ 模
$p$。}
\begin{enumerate}[resume]
  \item 证明\emph{威尔逊定理}：对质数 $p$，$(p-1)!
        \equiv -1 \pmod p$。\emph{（将每个剩余类与其
        逆元配对，哪些与自身配对？）}
  \item 设 $p$ 为奇质数，$m = \frac{p-1}2$。证明
        $(m!)^2 \equiv (-1)^{m+1} \pmod p$ \emph{（在 $(p-1)!$ 中，
        将每个因子 $k > m$ 替换为 $-(p - k)$）}。
  \item 由此得出：$-1$ 是模 $p$ 的二次剩余当且仅当 $p = 2$ 或 $p
        \equiv 1 \pmod 4$。\emph{（“仅当”部分：若 $x^2
        \equiv -1$，则 $x$ 在 $(\Z/p\Z)^\times$ 中的阶是多少，
        并由拉格朗日定理说明什么？）}
\end{enumerate}

\medskip
\noindent\textbf{第三部分 --- 分裂律。}
\begin{enumerate}[resume]
  \item 设 $p \equiv 1 \pmod 4$，且 $x$ 满足 $p \mid x^2 + 1 =
        (x + \iu)(x - \iu)$。证明 $p$ \emph{不}是
        高斯素数，并由第一部分推出 $p = a^2 + b^2$。
  \item 设 $p \equiv 3 \pmod 4$。直接证明 $p$ 不是
        两平方和\emph{（模 $4$ 的平方数）}，并推出
        $p$ 仍是高斯素数。
  \item 讨论 $p = 2$ 的情形：给出分解 $2 =
        -\iu(1+\iu)^2$，并验证 $1 + \iu$ 是高斯
        素数。（$2$ 是唯一的\emph{分歧}素数：可被
        高斯素数的平方整除，至多相差一个单位因子。）
  \item 综合高斯素数的分类，在相差
        单位的意义下：$1 + \iu$；整数 $p \equiv 3 \pmod 4$；
        共轭对 $\pi, \bar\pi$，范数为 $p \equiv 1 \pmod
        4$。在 $5 = (2+\iu)(2-\iu)$ 与 $3$ 上检验这一分类。
\end{enumerate}

\medskip
\noindent\textbf{第四部分 --- 两平方和定理。}
\begin{enumerate}[resume]
  \item 证明正向部分：若在分解式 $n =
        \prod p_i^{\alpha_i}$ 中每个 $\equiv 3 \pmod 4$ 的质数
        都以偶次幂出现，则 $n$ 是两个平方数之和。
        \emph{（婆罗摩笈多恒等式及第二--第三部分。）}
  \item 证明逆向部分：若 $n = a^2 + b^2 = N(a + \iu b)$，且
        $q \equiv 3 \pmod 4$ 整除 $n$，证明 $q$ 是
        高斯素数，整除 $a + \iu b$ 或 $a - \iu b$，
        并且它实际上同时整除 $a$ 与 $b$，再对 $n$ 用归纳法
        得出 $q$ 在 $n$ 中的指数为偶数。
  \item 叙述最终定理。在 $2025$、$2026$、$2027$
        中，哪些是两个平方数之和？\emph{（$2025 = 81 \cdot 25$；
        $2026 = 2 \cdot 1013$，$1013$ 为质数；$2027$ 为质数。）}
  \item （尾声）证明：$p \equiv 1 \pmod 4$ 的质数可以
        本质上唯一的方式表示为两个平方数之和：若 $p =
        a^2 + b^2 = c^2 + d^2$（正整数），则 $\{a, b\}
        = \{c, d\}$。\emph{（
        $\Z[\iu]$ 中分解的唯一性）}
\end{enumerate}

\medskip
\noindent\textbf{第五部分 --- 表示的计数：雅可比
公式与莱布尼茨级数。} 记 $r_2(n) = \#\{(a, b) \in
\Z^2 : a^2 + b^2 = n\}$（有序对，计入正负号与零），
并令 $\chi$ 为模 $4$ 的非平凡特征标：
$\chi(d) = +1$（若 $d \equiv 1$），$-1$（若 $d \equiv 3 \pmod4$），
$0$（若 $d$ 为偶数）。
\begin{enumerate}[resume]
  \item （热身，作为对比）哪些整数是
        两个平方数的\emph{差}？证明：$n = a^2 - b^2$
        对 $a, b \in \Z$ 成立，当且仅当 $n \not\equiv 2 \pmod 4$ ---
        无需环论，也不涉及可与下文相比拟的
        结构。
  \item 证明 $r_2(n)$ 是 $z \in \Z[\iu]$ 中
        满足 $N(z) = n$ 的元素个数。记 $n = 2^{a}\prod_jp_j^{b_j}
        \prod_kq_k^{c_k}$，其中 $p_j \equiv 1$，$q_k \equiv 3
        \pmod4$，利用第 11 题的分类与
        唯一分解，证明：这样的 $z$ 存在当且仅当所有
        $c_k$ 均为偶数，且在此情形下
        \[
        r_2(n) = 4\prod_j\,(b_j + 1) .
        \]
        \emph{（计数：$z = u\,(1+\iu)^{a}\prod_j\pi_j^{s_j}
        \bar\pi_j^{\,b_j - s_j}\prod_kq_k^{c_k/2}$，其中 $u$ 为
        单位，$0 \leq s_j \leq b_j$；这一列表为何
        既穷尽又无重复？）}
  \item 证明 $d \mapsto \chi(d)$ 是完全积性函数，
        由此推出 $n \mapsto
        \sum_{d \mid n}\chi(d)$ 是积性函数，并计算它
        在素数幂上的值：$1$（对 $2^a$）；$b + 1$（对
        $p^b$，$p \equiv 1$）；$1$ 或 $0$（对 $q^c$，$q \equiv 3$），
        视 $c$ 为偶数或奇数而定。
  \item 由此得出\emph{雅可比定理}：
        \[
        r_2(n) = 4\sum_{d \mid n}\chi(d) =
        4\bigl(d_1(n) - d_3(n)\bigr),
        \]
        其中 $d_i(n)$ 表示满足 $\equiv i \pmod 4$ 的因数个数。
        在 $n = 3, 5, 9, 25$ 上验证，并列出
        $16$ 种 $65$ 的表示。
  \item （圆）证明 $\sum_{n \leq x}r_2(n)$ 是
        $\Z^2$ 的格点中落在半径为
        $\sqrt x$ 的闭圆盘内的个数，并证明
        \[
        \sum_{n\leq x}r_2(n) = \pi x + O(\sqrt x)
        \]
        \emph{（每个格点对应一个单位正方形；比较
        面积，误差位于宽为
        $O(1)$ 的圆环中）}。
  \item （从算术角度读莱布尼茨级数）结合第 19--20 题：
        \[
        \sum_{d \leq x}\chi(d)\Bigl\lfloor\frac
        xd\Bigr\rfloor = \frac{\pi x}4 + O(\sqrt x),
        \]
        并推出 --- 小心去掉取整号 ---
        莱布尼茨级数
        \[
        1 - \frac13 + \frac15 - \frac17 + \dots = \frac\pi4 .
        \]
        奇数倒数的交错级数\emph{正是}
        $\equiv 1$ 的因数相对于
        $\equiv 3$ 的因数的平均超出量：这是借算术计算出的分析。
  \item （两平方和有多稀少？）证明：没有
        $\equiv 3 \pmod 4$ 的整数是两个平方数之和（有两种
        证法：模 $4$ 的平方数，或者
        第 17 题的奇偶判别法），故至少四分之一的整数
        被漏掉；并证明第 20 题中的平均值 $\frac1x\sum_{n\leq
        x}r_2(n) \to \pi$ 与之相容，
        密度为 $0$ 的可表示整数 ---
        举出表示方法异常多的整数
        （取许多个 $\equiv 1 \bmod 4$ 的素数之积）以
        说明一个趋于零的比例如何仍能承载
        一个正的平均值。（兰道证明了真正的密度以
        $1/\sqrt{\log x}$ 的速度衰减；这超出了我们的工具范围，但
        其机制现已清楚可见。）
\end{enumerate}

\medskip
\noindent\textbf{第六部分 --- 补充：本原
表示与毕达哥拉斯。}
\begin{enumerate}[resume]
  \item 称 $n = a^2 + b^2$ 的一个表示为\emph{本原}的，若
        $\gcd(a, b) = 1$。证明对 $n \geq 1$，它有
        本原表示当且仅当 $4 \nmid n$ 且没有素数
        $q \equiv 3 \pmod 4$ 整除 $n$。\emph{（对必要性，
        可沿用第 13 题的无穷递降法及
        模 $4$ 的平方数；对充分性，$z$ 仅由 $1 + \iu$
        及 $\pi_j$ 构成 --- 不用共轭 --- 并说明
        $a$ 与 $b$ 的公共素因子为何会迫使
        $\pi_j$ 与 $\bar\pi_j$，或 $(1+\iu)^2$，进入
        $z$。）}
  \item （毕达哥拉斯三元组）设 $a^2 + b^2 = c^2$，其中 $a, b,
        c$ 为正，$\gcd(a, b) = 1$ 且 $b$ 为偶数。证明
        $a + \iu b$ 与 $a - \iu b$ 在 $\Z[\iu]$ 中互素
        \emph{（若有公共的高斯素元，则它整除
        $2a$ 与 $2b$，而 $c$ 为奇数）}，由唯一分解性
        得 $a + \iu b = u(m + \iu n)^2$，
        其中 $u$ 为单位，并得出经典的参数化：
        在互换 $a$ 与 $b$ 的意义下，
        \[
        a = m^2 - n^2, \qquad b = 2mn, \qquad c = m^2 + n^2,
        \]
        其中 $m > n \geq 1$，且二者互素、奇偶性相反。
        恢复 $(3, 4, 5)$ 与 $(21, 20, 29)$，分别由 $(m, n) =
        (2, 1)$ 与 $(5, 2)$ 得到。
  \item （数值验证）取 $x = 25$。计算
        $r_2(n)$（$1 \leq n \leq 25$）由雅可比公式，
        验证非零值恰好出现在 $n = 1,
        2, 4, 5, 8, 9, 10, 13, 16, 17, 18, 20, 25$，并且
        \[
        \sum_{n \leq 25} r_2(n) = 80
        = 4\sum_{d \leq 25}\chi(d)
        \Bigl\lfloor\frac{25}d\Bigr\rfloor .
        \]
        验证半径为 $5$ 的闭圆盘含有 $81$
        个格点，并与 $\pi x \approx 78.5$ 比较：
        误差完全落在第
        20 题给出的 $O(\sqrt x)$ 界之内。
\end{enumerate}
\end{problem}
